\documentclass{../yalumba}

\usepackage{mathtools}
\usepackage{bm}

\title{Inherited Ecosystem Structure:\\A Spectral-Geometric Framework for\\Reference-Free Genomic Relatedness Detection}
\author{yalumba project}
\date{April 2026}
\reportslug{theoretical-framework-2026-04}

\setabstract{%
We propose that biological relatedness between individuals can be detected
by comparing the \emph{organizational structure} of their genomes rather than
their \emph{sequence content}.  In our framework, each genome is modeled as
a multi-scale ecosystem of interacting motif modules---groups of co-occurring
$k$-mers that form coherent functional or structural units within sequencing
reads.  We analyze this ecosystem via spectral geometry, constructing a
labeled interaction graph whose nodes are modules and whose edges encode
co-occurrence strength, then characterizing the graph through its Laplacian
spectrum, heat kernel signatures, and ecological role assignments.
Relatedness between two individuals is detected through \emph{structural
correspondence}: a continuous coherence field $\Phi: G_A \to G_B$ that maps
the ecosystem of one genome into the ecosystem of another, scored by how
well this mapping preserves local neighbourhood geometry.  This is
fundamentally different from identity-by-descent (IBD) methods, which
compare shared sequence tokens---identical $k$-mers, shared haplotype
segments, or concordant variant calls.  We have built a complete
from-scratch compute stack (12 packages, ${\sim}15{,}000$ lines of
TypeScript, custom GPU kernel pipeline) and tested more than 50 algorithms
across three datasets: a synthetic family trio, the GIAB Ashkenazi trio,
and the CEPH~1463 Platinum Pedigree (a three-generation, six-member
family that constitutes the hardest benchmark for reference-free methods).
Three key findings emerge from this systematic empirical program.
\textbf{First}, token-matching methods---from naive $k$-mer Jaccard through
sophisticated rare-run scanning---achieve strong separation on related versus
unrelated pairs but converge toward classical IBD proxies: the best
token method, rare-run P90, achieves $+583\%$ separation on CEPH~1463 and
is mathematically equivalent to reference-free IBD without alignment.
\textbf{Second}, structural methods---module persistence ($+3.36\%$),
coalition transfer ($+1.97\%$), spectral ecology ($+4.93\%$), and module
stability ($+4.90\%$ average but failing the weakest pair)---detect a
genuinely different signal, inherited organizational patterns, but this
signal is $100\times$ weaker than the rare-run signal and is confounded
by sequencing coverage variation.  \textbf{Third}, the gap between
structural and token methods has narrowed dramatically over six algorithm
generations in the spectral ecology family: from complete failure
($-26.87\%$ weakest gap in v1) through the spouse-confound breakthrough
(v4, spouses dropped from rank~1 to rank~8) to near-parity ($-0.61\%$
weakest gap in v7 via Sinkhorn-regularized transport with cooperative
stability weighting).  We identify three theoretical barriers to further
progress---coverage confounding, transport symmetry, and outlier
sensitivity---and propose solutions rooted in optimal transport
(Gromov--Wasserstein distance between ecosystem metric spaces), algebraic
topology (persistent homology of co-occurrence filtrations), and
differential geometry (gauge-theoretic field comparison via connection
curvature).  These three directions represent qualitative mathematical
leaps rather than incremental parameter tuning, and we argue that the
strongest case for the ecosystem approach will come not from beating
rare-run P90 at parent--child detection but from detecting relationships
or organizational structures that token-based methods fundamentally
cannot access.  This paper serves as a theoretical position statement for
external feedback: we present the complete intellectual framework, the
empirical trajectory, the open problems, and the proposed path forward.}

\begin{document}
\maketitle

\tableofcontents
\newpage


% ════════════════════════════════════════════════════════════════════════
% SECTION 1: INTRODUCTION
% ════════════════════════════════════════════════════════════════════════
\section{Introduction}
\label{sec:intro}

\subsection{The Reference-Free Relatedness Problem}
\label{sec:intro:problem}

Determining whether two individuals share recent common ancestors is a
foundational problem in human genetics.  The standard pipeline requires
multiple computationally expensive steps:

\begin{enumerate}
  \item \textbf{Alignment} to a reference genome (BWA-MEM, minimap2):
    hours per sample, reference-dependent.
  \item \textbf{Variant calling} (GATK HaplotypeCaller, DeepVariant):
    hours per sample, model-dependent.
  \item \textbf{Phasing} (SHAPEIT, Eagle): population-level, requires
    reference panels.
  \item \textbf{IBD detection} (GERMLINE~\cite{gusev2009germline},
    Refined IBD~\cite{browning2013refinedibd},
    IBIS~\cite{browning2020ibis}): segment matching on phased haplotypes.
\end{enumerate}

Each step introduces its own error modes and computational cost.
Alignment is reference-biased~\cite{liao2025sv}.  Variant calling has
systematic blind spots in repetitive regions.  Phasing requires large
reference panels and introduces switch errors.  The result is a pipeline
that works well for well-characterized populations with good reference
genomes, but degrades for under-represented populations, non-model
organisms, and metagenomic samples where no reference exists at all.

\begin{pullquote}
Can we detect relatedness directly from raw sequencing reads,
without alignment, variant calling, or phasing?
\end{pullquote}

This is the reference-free relatedness problem.  Prior work has shown
that $k$-mer sketching can estimate genomic distance
(Mash~\cite{ondov2016mash}), that alignment-free phylogenetics is
feasible~\cite{brinda2023phylogenetic}, and that reference-free
relatedness estimation is possible in principle~\cite{fan2021reference}.
Our contribution is a theoretical framework and systematic empirical
program that goes beyond $k$-mer counting to propose a fundamentally
different basis for comparison: the \emph{organizational structure} of
the genome as a multi-scale ecosystem.

\subsection{Two Paradigms for Genomic Comparison}
\label{sec:intro:paradigms}

We distinguish two paradigms that differ not in implementation details
but in what \emph{mathematical object} is compared between individuals.

\subsubsection{The Token Paradigm}

Genomes are bags of discrete tokens: $k$-mers, variants, haplotype
segments.  Relatedness is measured by shared tokens.  This is the
classical view underlying all IBD-based methods.

\begin{definition}
\textbf{Token comparison.}  Given two genomes $G_A$ and $G_B$
represented as sets of tokens $T_A$ and $T_B$, relatedness is a function
$f(T_A \cap T_B, T_A \cup T_B)$.  The Jaccard coefficient,
containment index, MinHash estimate, and IBD segment count are all
instances of this pattern.
\end{definition}

Token methods are intuitive, well-understood, and effective.  They have
a clear failure mode: they require tokens to be \emph{identical} across
individuals.  In genomic data, heterozygosity means that $k$-mers
within a shared IBD segment may differ at SNP positions.  At $k = 21$,
a single heterozygous SNP destroys $21$ $k$-mers.  With a SNP rate
of $\sim 0.1\%$ (one per $1{,}000$ bases), the probability that a
21-mer is identical between two haplotypes is approximately
$(1 - 0.001)^{21} \approx 0.979$.  This means ${\sim}2\%$ of $k$-mers
are disrupted by heterozygosity even within truly shared segments.

\subsubsection{The Ecosystem Paradigm (Our Proposal)}

Genomes are organized systems of interacting elements.  Relatedness is
measured by shared \emph{organizational structure}.  This does not
require any element to be identical across individuals---only that the
\emph{pattern of organization} is similar.

\begin{definition}
\textbf{Ecosystem comparison.}  Given two genomes $G_A$ and $G_B$
represented as labeled interaction graphs $\mathcal{G}_A = (V_A, E_A,
w_A, f_A)$ and $\mathcal{G}_B = (V_B, E_B, w_B, f_B)$, relatedness
is a function of the structural correspondence between $\mathcal{G}_A$
and $\mathcal{G}_B$: how well can the local neighbourhood geometry of
one be mapped to the other?
\end{definition}

The ecosystem paradigm is more robust to noise (a few disrupted $k$-mers
do not destroy the organizational structure) but harder to measure
(structural correspondence is a richer and more ambiguous quantity than
set intersection).  The central question of this research program is
whether the additional robustness translates into practical advantage.

\subsection{Symbiogenesis as Computational Metaphor}
\label{sec:intro:symbiogenesis}

Our framework draws on Lynn Margulis's theory of
symbiogenesis~\cite{margulis1998}: the idea that complex life arose not
through gradual mutation and selection alone, but through the
\emph{merging of simpler organisms into cooperative wholes}.
Mitochondria were once free-living bacteria; chloroplasts were once
cyanobacteria.  The eukaryotic cell is a consortium, not a monolith.

We apply this metaphor to genomic inheritance.  Each haplotype block
inherited from a parent is an ``organism'' carrying its internal
structure---its pattern of motif modules, their co-occurrence
relationships, their spectral geometry.  During meiosis, these organisms
recombine: maternal and paternal haplotype blocks interleave to form
a new consortium.  The child's genome is a \emph{symbiogenetic merger} of
parental organizational units.

\begin{keyfinding}
\textbf{The symbiogenesis metaphor predicts a specific testable
quantity}: related individuals should share not just individual genomic
elements but the \emph{cooperative structure} among those elements.
Two genomes that inherited the same haplotype block from a common
ancestor should exhibit the same local arrangement of motif modules---the
same co-occurrence edges, the same spectral local geometry, the same
ecological roles---within that block.
\end{keyfinding}

This prediction is what our coherence field framework attempts to test.

\subsection{Contributions and Scope}
\label{sec:intro:contributions}

This paper presents:

\begin{enumerate}
  \item A \textbf{theoretical framework} that models genomes as labeled
    interaction graphs and detects relatedness through spectral-geometric
    structural correspondence (Section~\ref{sec:framework}).
  \item A \textbf{comprehensive empirical survey} of 50+ algorithms
    across three datasets, tracing the evolution from naive $k$-mer
    overlap to sophisticated structural methods
    (Section~\ref{sec:trajectory}).
  \item A \textbf{formal analysis} of three theoretical barriers to
    further progress (Sections~\ref{sec:directions}
    and~\ref{sec:coverage}).
  \item \textbf{Three proposed mathematical directions}---optimal
    transport, persistent homology, and gauge-theoretic field
    comparison---each representing a qualitative leap beyond current
    methods (Section~\ref{sec:directions}).
  \item An \textbf{honest assessment} of what works, what fails, and
    where the ecosystem approach can or cannot compete with
    token methods (Section~\ref{sec:openquestions}).
\end{enumerate}

The computational infrastructure underlying this work is a
purpose-built monorepo (\texttt{yalumba}) containing 12 TypeScript
packages, a custom SPIR-V GPU compilation pipeline, and
${\sim}15{,}000$ lines of code written entirely from scratch with no
external bioinformatics libraries.


% ════════════════════════════════════════════════════════════════════════
% SECTION 2: THE EMPIRICAL TRAJECTORY
% ════════════════════════════════════════════════════════════════════════
\section{The Empirical Trajectory}
\label{sec:trajectory}

Before presenting theory, we describe the empirical program that
motivated it.  The trajectory from naive $k$-mer counting to structural
ecosystem comparison spans seven months, 50+ algorithm variants, and
three datasets of increasing difficulty.

\subsection{The Population Sharing Wall}
\label{sec:trajectory:wall}

The fundamental challenge for reference-free $k$-mer methods is that
most $k$-mers are shared between \emph{any} two same-population humans.
At $k = 21$, the probability that a randomly chosen 21-mer is identical
between two unrelated individuals is:

\begin{equation}
  P(\text{match}) = (1 - p_{\text{SNP}})^k
  \label{eq:pmatch}
\end{equation}

where $p_{\text{SNP}} \approx 0.001$ is the per-base SNP rate between
two haplotypes from the same population.  Substituting:

\begin{equation}
  P(\text{match}) = (1 - 0.001)^{21} = 0.999^{21} \approx 0.9792
  \label{eq:pmatch_val}
\end{equation}

This means approximately $97.9\%$ of 21-mers are shared between any two
individuals from the same population, regardless of kinship.  IBD
segments contribute an additional ${\sim}1{-}2\%$ of sharing on top of
this baseline.  The signal-to-noise ratio is $0.01$--$0.02$: a $1{-}2\%$
excess on a $98\%$ background.

\begin{limitation}
\textbf{The population sharing wall.}  Any method based on counting
shared $k$-mers must distinguish a ${\sim}1\%$ signal from a
${\sim}98\%$ background.  This is feasible with enough data (millions of
$k$-mers per sample) but creates enormous sensitivity to noise, coverage
variation, and sequencing error.
\end{limitation}

\subsection{Token-Based Methods: A Complete Survey}
\label{sec:trajectory:token}

We systematically tested every token-based comparison strategy we could
conceive.  Table~\ref{tab:token_methods} summarizes all families with
their best results on the CEPH~1463 benchmark (our hardest dataset:
three-generation pedigree with four parent--child relationships).

\begin{table}[H]
  \centering
  \tablestyle
  \caption{Token-based algorithm families tested, with best CEPH~1463
    results.  Separation = (mean related $-$ mean unrelated) / mean
    unrelated.  Gap = (weakest related $-$ strongest unrelated) / strongest
    unrelated.  A positive gap means \emph{all} related pairs score above
    \emph{all} unrelated pairs---perfect ordinal discrimination.}
  \label{tab:token_methods}
  \begin{tabular}{llcc}
    \toprule
    \rowcolor{table-header}
    \textcolor{white}{\textbf{Family}} &
    \textcolor{white}{\textbf{Best variant}} &
    \textcolor{white}{\textbf{Sep (\%)}} &
    \textcolor{white}{\textbf{Gap (\%)}} \\
    \midrule
    Set overlap & Jaccard, containment, MinHash & $<1$ & $<-50$ \\
    Run-length & Recombination distance & +12 & $-30$ \\
    Distribution & Bray--Curtis, Morisita--Horn & +3 & $-25$ \\
    Information & Entropy rate, NCD~\cite{li2004ncd} & +1.5 & $-40$ \\
    Multi-scale & Multi-scale run weighting & +180 & +50 \\
    Rare run-length & P90 with rarity filtering & \textbf{+583} & \textbf{+300} \\
    \bottomrule
  \end{tabular}
\end{table}

\ymarginnote{The rare-run P90 achieves $583\%$ separation---two orders
of magnitude above the baseline.}

The trajectory tells a clear story.  Naive methods (Jaccard, containment)
fail because of the population sharing wall: shared $k$-mers are
dominated by the universal $98\%$ background.  Run-length methods
improve by leveraging the spatial structure of sharing: IBD segments
create long \emph{runs} of consecutive shared $k$-mers, while
background sharing creates short, fragmented runs.  The P90 statistic
(90th percentile of run lengths) captures this efficiently.

The breakthrough came from combining run-length analysis with
\emph{rarity filtering}: restricting attention to $k$-mers that appear
in fewer than $50\%$ of samples.  This automatically removes the
population baseline, and the remaining rare $k$-mers are
overwhelmingly informative about kinship.

\subsection{The Rarity Insight}
\label{sec:trajectory:rarity}

The rare-run P90 method works as follows:

\begin{enumerate}
  \item Identify $k$-mers shared between a pair of samples.
  \item Filter to $k$-mers that appear in fewer than half the cohort
    (``rare'' $k$-mers).
  \item Compute the lengths of consecutive runs of shared rare $k$-mers
    along each read.
  \item Take the 90th percentile of run lengths as the similarity score.
\end{enumerate}

\begin{keyfinding}
\textbf{Rarity filtering is mathematically equivalent to automatic
variant detection.}  A $k$-mer that appears in fewer than $50\%$ of
samples from the same population is, with high probability, tagging a
variant (SNP, indel, or structural variant) that is not fixed in the
population.  The rare-run P90 is therefore detecting runs of shared
rare variants---which is precisely what IBD detectors do, but without
alignment or variant calling.
\end{keyfinding}

This insight reveals both the power and the limitation of the
approach.  The rare-run P90 achieves $+583\%$ separation because it is
essentially performing reference-free IBD detection.  But it is doing
\emph{exactly what GERMLINE does}, just without the alignment step.
It does not introduce any fundamentally new signal.

\subsection{The Convergence Problem}
\label{sec:trajectory:convergence}

As token methods improve, they converge toward classical IBD proxies:

\begin{itemize}
  \item Jaccard $\to$ MinHash $\to$ weighted MinHash $\approx$
    \emph{Mash}~\cite{ondov2016mash}
  \item Run-length $\to$ weighted runs $\to$ rare-filtered runs $\approx$
    \emph{reference-free GERMLINE}
  \item Distribution similarity $\to$ IDF-weighted cosine $\approx$
    \emph{TF-IDF}~\cite{salton1988tfidf} \emph{for genomes}
\end{itemize}

Each improvement makes the token method more closely approximate what a
traditional IBD detector computes.  This is \emph{not} a criticism---the
methods work precisely because IBD is the correct biological signal.
But it means that the token paradigm has a known ceiling: it cannot
detect anything that IBD-based methods cannot also detect (given
sufficient alignment and phasing quality).

\begin{pullquote}
The token paradigm converges toward classical IBD.\\
The question is whether the ecosystem paradigm can detect something else.
\end{pullquote}

\subsection{Structural Methods: A Different Signal}
\label{sec:trajectory:structural}

In parallel with the token program, we developed algorithms based on
the ecosystem paradigm.  These methods construct module-level
representations of each genome and compare their organizational
structure.  Table~\ref{tab:structural_methods} summarizes results.

\begin{table}[H]
  \centering
  \tablestyle
  \caption{Structural algorithm families on CEPH~1463.  All detect a
    genuinely different signal from token methods but achieve weaker
    separation.}
  \label{tab:structural_methods}
  \begin{tabular}{lccl}
    \toprule
    \rowcolor{table-header}
    \textcolor{white}{\textbf{Algorithm}} &
    \textcolor{white}{\textbf{Sep (\%)}} &
    \textcolor{white}{\textbf{Gap (\%)}} &
    \textcolor{white}{\textbf{Status}} \\
    \midrule
    Module Persistence v3 & +3.36 & +0.13 & \textbf{PASSES} \\
    Coalition Transfer v3 & +1.97 & +0.06 & \textbf{PASSES} \\
    Module Stability v4 & +4.90 & $-$1.04 & Fails weakest \\
    Spectral Ecology v3 & +4.93 & $-$54.67 & Fails \\
    Spectral Ecology v4 & +0.02 & $-$4.99 & Spouse broken \\
    Spectral Ecology v5 & $-$0.74 & $-$11.29 & Scale consist. \\
    Spectral Ecology v6 & n/a & $-$3.21 & Persistence \\
    Spectral Ecology v7 & n/a & $-$0.61 & Sinkhorn \\
    \bottomrule
  \end{tabular}
\end{table}

Module Persistence v3 and Coalition Transfer v3 achieve positive
weakest gap---meaning \emph{every} parent--child pair scores above
\emph{every} unrelated pair---but with separation orders of magnitude
smaller than the rare-run P90 ($+3.36\%$ vs.\ $+583\%$).

\begin{keyfinding}
\textbf{Structural methods detect a real but weak signal.}  The positive
gap of Module Persistence ($+0.13\%$) and Coalition Transfer ($+0.06\%$)
on CEPH~1463 demonstrates that organizational structure carries kinship
information independent of token identity.  But the signal-to-noise
ratio is ${\sim}100\times$ worse than token methods.
\end{keyfinding}

\subsection{The Spectral Ecology Arc: v1 through v7}
\label{sec:trajectory:arc}

The Spectral Ecology family represents our most sustained attempt to
extract structural signal through spectral-geometric methods.  Seven
versions trace a coherent intellectual trajectory, each version's failure
directly motivating the next.

\begin{table}[H]
  \centering
  \tablestyle
  \caption{The Spectral Ecology research trajectory.  Each version
    addresses a specific limitation of its predecessor.  Gap = weakest
    parent--child minus strongest non-parent--child, as percentage.}
  \label{tab:spectral_arc}
  \begin{tabular}{clcccl}
    \toprule
    \rowcolor{table-header}
    \textcolor{white}{\textbf{Ver}} &
    \textcolor{white}{\textbf{Paradigm}} &
    \textcolor{white}{\textbf{Sep}} &
    \textcolor{white}{\textbf{Gap}} &
    \textcolor{white}{\textbf{Sp.}} &
    \textcolor{white}{\textbf{Lesson learned}} \\
    \midrule
    v1 & Sequential adjacency
      & $-$0.7\% & $-$26.9\% & \#1
      & Graphs too sparse (2--5 edges) \\
    v2 & Within-sample spectra
      & +1.6\% & $-$22.5\% & \#1
      & Spectra have ceiling \\
    v3 & Cross-sample transport
      & +4.9\% & $-$54.7\% & \#1
      & Transport indiscriminate \\
    v4 & Single-scale coherence field
      & +0.02\% & $-$5.0\% & \#8
      & \textbf{Spouse confound broken} \\
    v5 & Multi-scale fields
      & $-$0.7\% & $-$11.3\% & \#6
      & Scale consistency real \\
    v6 & Persistence fields
      & n/a & $-$3.2\% & mid
      & Persistence narrows gap \\
    v7 & Sinkhorn + cooperative
      & n/a & $-$0.6\% & mid
      & \textbf{Near parity} \\
    \bottomrule
  \end{tabular}
\end{table}

\ymarginnote{Seven versions. The weakest gap improved from $-26.9\%$ to
$-0.6\%$---a $44\times$ reduction.}

The critical transition occurs at v4, when the \emph{coherence field}
is introduced.  Versions v1--v3 attempt to compare ecosystem
\emph{summaries} (graph spectra, heat kernel traces, bipartite
transport scores).  These summaries capture some structural information
but cannot distinguish spouses (who have similar global graph structure
due to similar coverage) from parent--child pairs (who share specific
local neighbourhood arrangements due to IBD).

Version v4 introduces a continuous field $\Phi: G_A \to G_B$ that
maps every module in sample $A$ to a position in sample $B$'s ecosystem
space, scored by whether the mapping preserves local neighbourhood
geometry.  This is the first method to break the spouse confound:
spouses have globally similar ecosystems but locally dissimilar
neighbourhoods, producing incoherent fields with high distortion.

Versions v5--v7 refine the field framework: multi-scale consistency (v5),
topological persistence (v6), and Sinkhorn-regularized optimal transport
with cooperative stability weighting (v7).  The weakest gap contracts
from $-26.87\%$ (v1) to $-0.61\%$ (v7), a $44\times$ improvement that
demonstrates convergence toward the passing threshold.


% ════════════════════════════════════════════════════════════════════════
% SECTION 3: THEORETICAL FRAMEWORK
% ════════════════════════════════════════════════════════════════════════
\section{Theoretical Framework}
\label{sec:framework}

We now formalize the ecosystem comparison approach.  The framework
proceeds in four stages: genome representation, spectral analysis,
ecological characterization, and structural correspondence.

\subsection{Genome as Labeled Interaction Graph}
\label{sec:framework:graph}

\begin{definition}
\textbf{Genomic ecosystem graph.}  Given a set of sequencing reads from
an individual, we construct a labeled interaction graph
\[
  \mathcal{G} = (V, E, w, f)
\]
where:
\begin{itemize}
  \item $V = \{m_1, \ldots, m_n\}$ is a set of \emph{modules}---clusters
    of co-occurring $k$-mers identified by community detection on the
    $k$-mer co-occurrence graph;
  \item $E \subseteq V \times V$ is a set of \emph{co-occurrence edges}
    connecting modules that appear together in sequencing reads more
    often than expected under independence;
  \item $w: E \to \mathbb{R}_{>0}$ assigns each edge a weight
    proportional to the strength of co-occurrence;
  \item $f: V \to \mathbb{R}^d$ assigns each module a $d$-dimensional
    feature vector encoding its local structural properties (support
    count, diversity, connectivity, spectral embedding coordinates).
\end{itemize}
\end{definition}

The key departure from token-based representations is that $\mathcal{G}$
encodes \emph{relationships between modules}, not just the modules
themselves.  Two genomes can share the same set of modules (high token
overlap) but arrange them differently (low structural correspondence),
or vice versa.

Module construction follows a three-stage pipeline:

\begin{enumerate}
  \item \textbf{Motif extraction.}  Extract all canonical $k$-mers
    (lexicographically smaller of forward/reverse complement) from reads
    using a Rabin--Karp rolling hash~\cite{karp1987rabin}.
  \item \textbf{Co-occurrence graph.}  Build a weighted graph where
    nodes are frequent $k$-mers and edges connect $k$-mers that
    co-occur within a genomic window (typically 50--150\,bp).
  \item \textbf{Community detection.}  Identify densely connected
    subgraphs as modules using cohesion-based clustering (minimum support
    and minimum internal edge density thresholds).
\end{enumerate}

In our implementation, typical parameters are $k = 15$, window size
50--150\,bp, minimum support 3--5, and minimum cohesion 0.25--0.30.
This produces 300--700 modules per sample in the CEPH~1463 dataset,
with 200--1{,}400 co-occurrence edges depending on sequencing depth.

\subsection{The Laplacian and Its Spectrum}
\label{sec:framework:laplacian}

The organizational structure of $\mathcal{G}$ is analyzed through its
graph Laplacian.  Let $A \in \mathbb{R}^{n \times n}$ be the weighted
adjacency matrix and $D = \mathrm{diag}(d_1, \ldots, d_n)$ the degree
matrix where $d_i = \sum_j A_{ij}$.

\begin{definition}
\textbf{Normalized Laplacian.}  The symmetric normalized Laplacian is
\begin{equation}
  L_{\mathrm{sym}} = I - D^{-1/2} A D^{-1/2}
  \label{eq:laplacian}
\end{equation}
with eigendecomposition $L_{\mathrm{sym}} = \Phi \Lambda \Phi^T$, where
$\Lambda = \mathrm{diag}(\lambda_0, \lambda_1, \ldots, \lambda_{n-1})$
with $0 = \lambda_0 \leq \lambda_1 \leq \cdots \leq \lambda_{n-1}
\leq 2$.
\end{definition}

The spectrum $\{\lambda_i\}$ encodes global structural properties:

\begin{itemize}
  \item $\lambda_0 = 0$ always (constant eigenvector; connected
    components).
  \item $\lambda_1$ = \emph{algebraic connectivity} (Fiedler
    value~\cite{chung1997spectral}): measures how well-connected
    the graph is.  Small $\lambda_1$ indicates a bottleneck or
    near-disconnection.
  \item The multiplicity of $\lambda_0 = 0$ equals the number of
    connected components.
  \item The spectral gap $\lambda_1 - \lambda_0 = \lambda_1$ controls
    the mixing time of random walks on the graph.
\end{itemize}

For ecosystem graphs with 300--700 nodes, we compute the full spectrum
via Jacobi eigendecomposition in sub-second time.  No iterative
approximations are needed.

The \emph{random walk normalized Laplacian} $L_{\mathrm{rw}} = I -
D^{-1}A$ is related by $L_{\mathrm{rw}} = D^{1/2} L_{\mathrm{sym}}
D^{-1/2}$ and shares the same eigenvalues.  We use $L_{\mathrm{sym}}$
for eigendecomposition (since it is symmetric and amenable to
standard solvers) and interpret results through the random walk
perspective.

\subsection{Heat Kernel Signature as Ecosystem Fingerprint}
\label{sec:framework:hks}

The heat kernel on $\mathcal{G}$ describes how a unit of ``heat''
placed at a node diffuses through the graph over time.  It is defined
as the matrix exponential of the Laplacian:

\begin{equation}
  K_t = \exp(-t L_{\mathrm{sym}})
  = \sum_{k=0}^{n-1} e^{-t \lambda_k} \, \phi_k \phi_k^T
  \label{eq:heat_kernel}
\end{equation}

where $\phi_k$ is the $k$-th eigenvector.  The \emph{heat kernel
signature} (HKS) of node $v$ at timescale $t$ is:

\begin{equation}
  \mathrm{HKS}(v, t) = K_t(v, v) = \sum_{k=0}^{n-1}
  e^{-t \lambda_k} \, \phi_k(v)^2
  \label{eq:hks}
\end{equation}

\begin{keyfinding}
\textbf{The HKS is a multi-scale structural descriptor.}  At small $t$,
HKS captures the immediate neighbourhood of $v$ (how many edges, how
dense).  At large $t$, HKS captures $v$'s global position in the graph
(central versus peripheral).  The full signature $\mathrm{HKS}(v,
\cdot): \mathbb{R}_{>0} \to \mathbb{R}$ is a complete characterization
of $v$'s structural role up to isospectral ambiguity.
\end{keyfinding}

We evaluate the HKS at 10 logarithmically spaced timescales from $t =
0.1$ to $t = 100$, producing a 10-dimensional descriptor per module.
The graph-level trace $\sum_v \mathrm{HKS}(v, t) = \mathrm{trace}(K_t)
= \sum_k e^{-t\lambda_k}$ is a permutation-invariant graph fingerprint
used in our within-sample spectral comparison (Spectral Ecology v2).

\subsection{Ecological Role Assignment via Spectral Embedding}
\label{sec:framework:roles}

The first $d$ eigenvectors of $L_{\mathrm{sym}}$ define a spectral
embedding $\psi: V \to \mathbb{R}^d$ where $\psi(v) = (\phi_1(v),
\ldots, \phi_d(v))$.  This embedding maps structurally similar nodes to
nearby points in $\mathbb{R}^d$, regardless of their identity (which
$k$-mers they contain).

We classify modules into ecological roles based on their spectral
embedding coordinates and graph-theoretic properties:

\begin{table}[H]
  \centering
  \tablestyle
  \caption{Ecological role taxonomy for genomic modules.  Classification
    is based on spectral embedding coordinates and local graph
    properties.}
  \label{tab:roles}
  \begin{tabular}{lll}
    \toprule
    \rowcolor{table-header}
    \textcolor{white}{\textbf{Role}} &
    \textcolor{white}{\textbf{Graph property}} &
    \textcolor{white}{\textbf{Biological interpretation}} \\
    \midrule
    Core & High degree, central embedding & Conserved genomic structure \\
    Satellite & Low degree, peripheral embedding & Variant-rich region \\
    Bridge & High betweenness, mid embedding & Connects genomic domains \\
    Scaffold & High clustering coefficient & Dense local structure \\
    Boundary & Spectral gap position & Domain transition zone \\
    \bottomrule
  \end{tabular}
\end{table}

The role distribution of a genome---e.g., ``8 cores, 85 satellites, 356
bridges, 12 scaffolds, 4 boundaries''---is a coarse structural
fingerprint that is independent of which $k$-mers populate each module.
Two genomes can have identical role distributions despite sharing no
$k$-mers.

\subsection{Coherence Fields: The Key Theoretical Object}
\label{sec:framework:fields}

The central theoretical contribution of our framework is the
\emph{coherence field}: a continuous mapping between ecosystem spaces
that measures structural correspondence.

\begin{definition}
\textbf{Coherence field.}  Given two genomic ecosystem graphs
$\mathcal{G}_A$ and $\mathcal{G}_B$ with patch tensor representations
$\{t_i^A\}_{i=1}^{n_A} \subset \mathbb{R}^d$ and $\{t_j^B\}_{j=1}^{n_B}
\subset \mathbb{R}^d$, the coherence field $\Phi: \mathcal{G}_A \to
\mathbb{R}^d$ maps each module in $A$ to a position in $B$'s ecosystem
space via kernel regression:
\begin{equation}
  \Phi(m_i^A) = \frac{\sum_{j=1}^{n_B} K(t_i^A, t_j^B) \cdot t_j^B}
                      {\sum_{j=1}^{n_B} K(t_i^A, t_j^B)}
  \label{eq:coherence_field}
\end{equation}
where $K(x, y) = \exp\!\left(-\|x - y\|^2 / 2\sigma^2\right)$ is a
Gaussian kernel with auto-scaled bandwidth $\sigma$ equal to the median
pairwise distance.
\end{definition}

The coherence field is not a hard assignment (each $A$-module maps to
one $B$-module) but a soft, continuous mapping.  Each $A$-module maps to
a \emph{weighted blend} of $B$-module positions, with weights
determined by structural similarity in patch-tensor space.

\ymarginnote{The coherence field is the ecosystem analogue of sequence
alignment.  Alignment maps bases to bases; the field maps ecological
neighbourhoods to ecological neighbourhoods.}

The patch tensor $t_i \in \mathbb{R}^d$ (with $d = 13$ in our
implementation) encodes the local structural environment of each module:

\begin{itemize}
  \item Local Laplacian eigenvalues $\lambda_1, \ldots, \lambda_6$ of
    the subgraph induced by the module's 1-hop neighbourhood: the
    ``spectral shape'' of the local ecosystem.
  \item Role composition (5 fractions): the proportion of core,
    satellite, bridge, scaffold, and boundary modules in the
    neighbourhood.
  \item Patch density: fraction of possible edges present in the local
    subgraph.
  \item Relative patch size: number of nodes in the patch divided by
    total graph size.
\end{itemize}

\subsection{Cooperative Stability: Measuring Field Quality}
\label{sec:framework:stability}

The coherence field alone does not produce a similarity score.  We need
to measure how \emph{well-behaved} the field is---whether it represents
a coherent structural correspondence or a random, meaningless mapping.

\begin{definition}
\textbf{Cooperative stability.}  The cooperative stability of a coherence
field $\Phi: \mathcal{G}_A \to \mathbb{R}^d$ is the smoothness of the
mapping with respect to the graph structure of $A$:
\begin{equation}
  S(\Phi) = \exp\!\left(-\mathrm{CV}\!\left(\left\{
    \|\Phi(m_i) - \Phi(m_j)\| : (m_i, m_j) \in E_A
  \right\}\right)\right)
  \label{eq:coop_stability}
\end{equation}
where $\mathrm{CV}$ denotes the coefficient of variation (standard
deviation divided by mean) of the set of edge-mapped distances.
\end{definition}

High cooperative stability ($S \to 1$) means that neighbouring modules
in $A$ map to nearby positions in $B$'s ecosystem space---the field
preserves local neighbourhood structure.  Low cooperative stability
($S \to 0$) means the mapping is erratic: neighbours in $A$ map to
distant, unrelated positions in $B$.

\begin{keyfinding}
\textbf{What cooperative stability actually measures.}
Consider modules $m_i$ and $m_j$ that are neighbours in genome $A$
(they co-occur in reads, their ecological niches overlap).  Cooperative
stability asks: ``do the corresponding regions of genome $B$ also
place these modules' analogues in neighbouring positions?''

This is NOT asking ``do $A$ and $B$ share the same modules'' (token
matching).  It is asking ``do $A$ and $B$ \emph{organize} their modules
in the same way?''
\end{keyfinding}

\subsection{Why Cooperative Stability Should Detect Relatedness}
\label{sec:framework:why}

The theoretical argument for cooperative stability as a relatedness
signal proceeds as follows:

\begin{enumerate}
  \item \textbf{IBD segments are inherited intact.}  When a child
    inherits a haplotype block from a parent, the entire arrangement of
    motifs, their co-occurrence patterns, and their local spectral
    geometry within that block is inherited intact (modulo sequencing
    errors and somatic mutations).

  \item \textbf{Within IBD segments, arrangement is identical.}  The
    coherence field restricted to modules from an IBD segment should be
    a near-identity mapping: each $A$-module maps to the ``same'' module
    in $B$.  This produces perfect cooperative stability within the
    segment.

  \item \textbf{Between non-IBD segments, arrangement is random.}  For
    genomic regions that are \emph{not} IBD, $A$ and $B$ inherited
    different haplotype blocks from different ancestors.  Even if
    individual motifs happen to match (due to population sharing),
    their arrangements are independent.  The coherence field over these
    regions has no reason to be smooth.

  \item \textbf{Global cooperative stability reflects IBD fraction.}
    Parent--child pairs share $\sim 50\%$ of their genome IBD.  The
    coherence field is smooth over the IBD half and erratic over the
    non-IBD half.  The aggregate smoothness should be higher than for
    unrelated pairs (where IBD fraction is near zero).
\end{enumerate}

\begin{limitation}
\textbf{Why cooperative stability currently does not fully work.}
The argument above assumes that we can distinguish ``smooth over IBD
segments'' from ``smooth by coincidence.''  In practice, three barriers
prevent this (detailed in Section~\ref{sec:coverage}):
\begin{enumerate}
  \item Coverage confound: similar sequencing depth creates similar
    module counts and graph densities, producing spurious smoothness.
  \item Symmetry: computing only $\Phi_{A \to B}$ (not also $\Phi_{B
    \to A}$) allows directional artifacts.
  \item Outlier sensitivity: atypical samples with unusual role
    distributions (e.g., NA12891 with 356 bridge modules) produce
    distinctive patch tensors that create artificially consistent
    fields.
\end{enumerate}
\end{limitation}


% ════════════════════════════════════════════════════════════════════════
% SECTION 4: THREE PROPOSED THEORETICAL DIRECTIONS
% ════════════════════════════════════════════════════════════════════════
\section{Three Proposed Theoretical Directions}
\label{sec:directions}

The coherence field framework (Spectral Ecology v4--v7) has narrowed
the gap from $-26.87\%$ to $-0.61\%$ but appears to be approaching
diminishing returns within the kernel-regression paradigm.  We propose
three mathematically distinct directions that represent qualitative leaps
rather than incremental refinements.

\subsection{Gromov--Wasserstein Distance Between Ecosystem Metric Spaces}
\label{sec:directions:gw}

\subsubsection{Motivation}

The coherence field $\Phi$ constructs an explicit point-to-point mapping
between modules.  This is sensitive to outliers and requires choosing a
kernel bandwidth.  An alternative is to compare the \emph{metric spaces}
directly, without any correspondence.

\begin{definition}
\textbf{Gromov--Wasserstein distance~\cite{memoli2011gromov}.}  Given two
metric measure spaces $(X, d_X, \mu_X)$ and $(Y, d_Y, \mu_Y)$, the
Gromov--Wasserstein distance is
\begin{equation}
  \mathrm{GW}(X, Y) = \inf_{\pi \in \Pi(\mu_X, \mu_Y)}
  \left(
    \iint_{X \times Y} \iint_{X \times Y}
    |d_X(x, x') - d_Y(y, y')|^2 \,
    d\pi(x,y) \, d\pi(x', y')
  \right)^{1/2}
  \label{eq:gw}
\end{equation}
where $\Pi(\mu_X, \mu_Y)$ is the set of couplings (transport plans)
between $\mu_X$ and $\mu_Y$.
\end{definition}

In words: GW distance finds the optimal soft correspondence between $X$
and $Y$ that minimizes the distortion of \emph{internal distances}.  It
asks: ``can the internal geometry of $X$ be isometrically mapped to
$Y$?''

\subsubsection{Application to Genomic Ecosystems}

For genomes $A$ and $B$:
\begin{itemize}
  \item $X = \{t_1^A, \ldots, t_{n_A}^A\} \subset \mathbb{R}^{13}$:
    patch tensors of $A$'s modules.
  \item $Y = \{t_1^B, \ldots, t_{n_B}^B\} \subset \mathbb{R}^{13}$:
    patch tensors of $B$'s modules.
  \item $d_X(i, i') = \|t_i^A - t_{i'}^A\|$: Euclidean distance in
    patch-tensor space.
  \item $\mu_X, \mu_Y$: uniform measures (all modules equally weighted)
    or support-weighted measures.
\end{itemize}

\begin{keyfinding}
\textbf{Why GW might work.}  GW distance compares the \emph{shape} of
the internal distance distribution, not the positions of individual
modules.  Related individuals share inherited arrangement
patterns---modules that are close together in one genome should be close
together in the other---producing similar internal distance
distributions.  Unrelated individuals share random arrangements,
producing different distributions.

Critically, GW is \textbf{inherently size-invariant}: it does not
require the two ecosystems to have the same number of modules.  This
directly addresses the coverage confound, where different sequencing
depths produce different module counts.
\end{keyfinding}

\subsubsection{Computational Feasibility}

The exact GW problem is NP-hard in general, but efficient algorithms
exist for the entropic-regularized version.  The Sinkhorn-based solver
of Peyré et al.~\cite{peyre2019computational} computes
$\mathrm{GW}_\varepsilon$ in $O(n^2 \log n)$ time per outer iteration,
with $n = \max(n_A, n_B)$.  For ecosystem graphs with 400--700 modules,
this is feasible in seconds.

\subsubsection{Comparison to Coherence Fields}

\begin{table}[H]
  \centering
  \tablestyle
  \caption{Coherence fields vs.\ Gromov--Wasserstein distance.}
  \label{tab:gw_comparison}
  \begin{tabular}{lll}
    \toprule
    \rowcolor{table-header}
    \textcolor{white}{\textbf{Property}} &
    \textcolor{white}{\textbf{Coherence field}} &
    \textcolor{white}{\textbf{GW distance}} \\
    \midrule
    Correspondence & Soft (kernel regression) & Optimal (transport plan) \\
    Symmetry & Asymmetric ($A \to B$) & Symmetric by construction \\
    Size invariance & No (module count matters) & Yes (metric space) \\
    Outlier sensitivity & High (kernel dominated) & Moderate (transport) \\
    Metric? & No & Yes (a proper distance) \\
    Complexity & $O(n_A \cdot n_B)$ & $O(n^2 \log n)$ per iteration \\
    \bottomrule
  \end{tabular}
\end{table}

\subsection{Persistent Homology of Co-Occurrence Complexes}
\label{sec:directions:ph}

\subsubsection{Motivation}

Spectral methods capture the ``shape'' of the ecosystem graph through
eigenvalues, but they are blind to certain topological features.  Two
graphs can have identical spectra but different topology (co-spectral
non-isomorphic graphs).  Persistent
homology~\cite{edelsbrunner2010computational} provides a complementary
topological characterization that captures features invisible to the
spectrum.

\subsubsection{Filtration Construction}

Starting from the weighted co-occurrence graph, we construct a
filtration of simplicial complexes:

\begin{enumerate}
  \item \textbf{Threshold filtration.}  At threshold $\varepsilon$,
    include all edges with weight $\geq \varepsilon$ and all cliques
    (complete subgraphs) induced by those edges.  As $\varepsilon$
    decreases from $\max(w)$ to $0$, more edges and higher-dimensional
    simplices appear.

  \item \textbf{Vietoris--Rips filtration.}  Embed modules in
    patch-tensor space and build the Rips complex at increasing
    radius.  This captures metric-space topology rather than graph
    topology.
\end{enumerate}

\begin{definition}
\textbf{Persistence diagram.}  For each topological feature (connected
component, loop, void, \ldots), the persistence diagram records a
\emph{birth time} $b$ (the threshold at which the feature first appears)
and a \emph{death time} $d$ (the threshold at which the feature is
destroyed by merger or filling).  The persistence $d - b$ measures the
significance of the feature: long-lived features represent robust
structural properties, while short-lived features are noise.
\end{definition}

\subsubsection{What Persistent Homology Captures}

\begin{itemize}
  \item $H_0$ (\textbf{connected components}): how the ecosystem
    fragments into clusters as edge weight threshold increases.
    Long-lived $H_0$ features correspond to robust subcommunities of
    modules.
  \item $H_1$ (\textbf{loops}): cycles in the co-occurrence graph that
    persist across thresholds.  These may indicate regulatory circuits,
    repeat structures, or circular dependencies between modules.
  \item $H_2$ (\textbf{voids}): higher-dimensional cavities in the
    simplicial complex.  Rare in genomic data but potentially informative
    about mesoscale organization.
\end{itemize}

\begin{keyfinding}
\textbf{Why persistent homology might work.}  Persistence diagrams are
\emph{stable} under small perturbations (the Bottleneck Stability
Theorem~\cite{edelsbrunner2010computational}): if two ecosystems differ
by small noise, their persistence diagrams are close.  They are
\emph{inherently size-invariant}: the number of features depends on
topology, not on the number of nodes.  And they capture
\emph{complementary information} to spectral methods: two ecosystems
with identical spectra can have different persistence diagrams.
\end{keyfinding}

\subsubsection{Comparison Between Genomes}

Given persistence diagrams $D_A$ and $D_B$ for two genomes, we compute
their distance using the $p$-Wasserstein distance between diagrams:

\begin{equation}
  W_p(D_A, D_B) = \left(\inf_{\gamma: D_A \to D_B}
  \sum_{x \in D_A} \|x - \gamma(x)\|_\infty^p \right)^{1/p}
  \label{eq:wasserstein_diagram}
\end{equation}

where the infimum is over all bijections $\gamma$ between the (augmented)
diagrams.  The augmentation allows unmatched points to be paired with
their projection onto the diagonal (birth = death), representing
ephemeral features.

\subsubsection{Computational Cost}

Computing the persistence diagram from a weighted graph on $n$ nodes
requires $O(n^3)$ time for the boundary matrix reduction (or
$O(n^\omega)$ using fast matrix multiplication).  For $n = 700$ modules,
this is sub-second.  The Wasserstein distance between two diagrams with
$m$ points is $O(m^3)$ via the Hungarian algorithm, or
$O(m^{1+\varepsilon})$ using recent geometric
algorithms~\cite{kerber2017geometry}.

\subsection{Gauge-Theoretic Field Comparison}
\label{sec:directions:gauge}

\subsubsection{Motivation}

The coherence field $\Phi$ is scored by its pointwise smoothness
(cooperative stability) and its pointwise accuracy (curvature mismatch,
distortion).  Both measures are sensitive to outliers: a few modules
with atypical patch tensors can dominate the aggregate score.

We propose reframing the field comparison in terms of \emph{gauge
theory}: instead of comparing $\Phi$ pointwise, we compare the
\emph{connection}---the way $\Phi$ varies along edges.

\subsubsection{The Field as a Fiber Bundle Section}

\begin{definition}
\textbf{Coherence bundle.}  Let $\mathcal{G}_A = (V_A, E_A)$ be the
base graph.  Over each node $v \in V_A$, attach a fiber
$F_v \cong \mathbb{R}^d$ representing the possible positions in $B$'s
ecosystem space.  The coherence field $\Phi$ is a section of this
bundle: $\Phi(v) \in F_v$ for each $v$.
\end{definition}

The \emph{connection} on this bundle is defined by the transition
functions along edges:

\begin{equation}
  \nabla_{ij} \Phi = \Phi(m_j) - P_{ij}(\Phi(m_i))
  \label{eq:connection}
\end{equation}

where $P_{ij}: F_i \to F_j$ is the parallel transport operator (in
the simplest case, the identity; in a more sophisticated version,
the optimal rotation/scaling from the tangent space at $i$ to the
tangent space at $j$).

\subsubsection{Curvature as a Relatedness Invariant}

The \emph{curvature} of the connection measures the holonomy around
closed loops in the graph.  For a triangle $(i, j, k)$ in
$\mathcal{G}_A$:

\begin{equation}
  R_{ijk} = \nabla_{jk} + \nabla_{ki} + \nabla_{ij}
  = \Phi(k) - P_{jk}(\Phi(j)) + \Phi(i) - P_{ki}(\Phi(k))
    + \Phi(j) - P_{ij}(\Phi(i))
  \label{eq:curvature}
\end{equation}

If the parallel transport is trivial ($P = \mathrm{id}$), this simplifies
to the ``boundary'' of $\Phi$ around the triangle.

\begin{keyfinding}
\textbf{Why curvature should detect relatedness.}  Related individuals
have coherence fields with \emph{low curvature} (flat connections):
the inherited organizational structure induces a \emph{consistent}
mapping that ``commutes'' around loops.  If modules $i \to j \to k$ form
a triangle in $A$'s ecosystem, and their images $\Phi(i), \Phi(j),
\Phi(k)$ form a corresponding triangle in $B$'s ecosystem space, then
the holonomy around the loop is zero.  Unrelated individuals produce
\emph{high-curvature} fields: the mapping is locally reasonable
(each edge maps smoothly) but globally inconsistent (going around
a loop returns to a different point).
\end{keyfinding}

\subsubsection{Aggregate Curvature Score}

The total curvature is aggregated over all triangles in $\mathcal{G}_A$:

\begin{equation}
  \mathcal{R}(\Phi) = \frac{1}{|T(\mathcal{G}_A)|}
  \sum_{(i,j,k) \in T(\mathcal{G}_A)} \|R_{ijk}\|^2
  \label{eq:total_curvature}
\end{equation}

where $T(\mathcal{G}_A)$ is the set of triangles in $A$'s graph.  Low
$\mathcal{R}$ indicates a flat (consistent) field; high $\mathcal{R}$
indicates an inconsistent field.

The relatedness score is then $S_{\mathrm{gauge}} = \exp(-\alpha \cdot
\mathcal{R}(\Phi))$ for an appropriate scaling constant $\alpha$.

\begin{limitation}
\textbf{This is the most speculative direction.}  Gauge-theoretic
methods have been applied to image analysis (shape matching via
connections on fiber bundles~\cite{nakahara2003geometry}) and to network
analysis (Ollivier--Ricci curvature on graphs), but not to genomic
ecosystem comparison.  The key open question is whether genomic
ecosystem graphs contain enough triangles to make curvature estimation
statistically meaningful.  In our CEPH~1463 data, typical ecosystem
graphs have 300--700 modules and 200--1{,}400 edges, yielding
$O(10^2)$--$O(10^3)$ triangles---likely sufficient, but empirical
validation is needed.
\end{limitation}

\subsubsection{Comparison of the Three Directions}

\begin{table}[H]
  \centering
  \tablestyle
  \caption{Summary of three proposed theoretical directions.}
  \label{tab:three_directions}
  \begin{tabular}{lcccl}
    \toprule
    \rowcolor{table-header}
    \textcolor{white}{\textbf{Direction}} &
    \textcolor{white}{\textbf{Size-inv.}} &
    \textcolor{white}{\textbf{Outlier}} &
    \textcolor{white}{\textbf{Metric?}} &
    \textcolor{white}{\textbf{Theory basis}} \\
    \midrule
    GW distance & Yes & Moderate & Yes & Optimal transport \\
    Persistent homology & Yes & Low & Yes & Algebraic topology \\
    Gauge curvature & Partial & Low & No & Differential geometry \\
    \bottomrule
  \end{tabular}
\end{table}

We believe these directions are complementary: GW distance captures
global metric structure, persistent homology captures topological
invariants, and gauge curvature captures local consistency.  A combined
score incorporating all three would leverage multiple mathematical
perspectives on the same underlying question: ``do these two ecosystems
have the same organizational structure?''


% ════════════════════════════════════════════════════════════════════════
% SECTION 5: THE COVERAGE CONFOUND
% ════════════════════════════════════════════════════════════════════════
\section{The Coverage Confound: A Formal Analysis}
\label{sec:coverage}

The single largest obstacle to structural relatedness detection is
confounding by sequencing coverage.  This section provides a formal
analysis of why coverage creates spurious similarity and proposes
four deconfounding strategies.

\subsection{Why Coverage Creates Spurious Similarity}
\label{sec:coverage:why}

\begin{definition}
\textbf{Null model.}  Consider two genomes with \emph{identical}
$k$-mer frequency distributions but \emph{independent} sequencing.
Let $c_A$ and $c_B$ denote their sequencing coverage (average number
of reads covering each base).
\end{definition}

Under this null model, the ecosystem graph properties scale
predictably with coverage:

\begin{enumerate}
  \item \textbf{Module count.}  A $k$-mer becomes a ``module motif'' when
    it exceeds a support threshold $\tau$.  The probability that a
    motif with population frequency $p$ exceeds threshold $\tau$ at
    coverage $c$ is approximately $P(\mathrm{Poisson}(p \cdot c) \geq
    \tau)$.  Thus module count $n \sim F(c)$ where $F$ is a monotone
    increasing function of coverage.

  \item \textbf{Edge count.}  Two motifs co-occur in a window when both
    are sequenced.  If motifs have independent coverage, the expected
    co-occurrence count scales as $O(c^2)$, and the number of edges
    exceeding a co-occurrence threshold scales as $O(c^2)$ as well.

  \item \textbf{Spectral gap.}  For Erd\H{o}s--R\'enyi random graphs
    $G(n, p)$ with edge probability $p \propto c^2/n^2$, the spectral gap
    $\lambda_1 \to 1$ as $np \to \infty$.  Higher coverage produces
    denser graphs with larger spectral gaps.

  \item \textbf{Heat kernel trace.}  The normalized trace
    $\mathrm{trace}(K_t) / n = \sum_k e^{-t\lambda_k} / n$ depends on
    the eigenvalue distribution, which itself depends on coverage through
    graph density.
\end{enumerate}

\begin{keyfinding}
\textbf{All standard spectral features correlate with coverage.}  Under
the null model, samples with similar coverage produce similar module
counts, edge counts, spectral gaps, and heat kernel traces---regardless
of kinship.  This explains why spouses (who often have similar
sequencing depth due to being processed together) are difficult to
separate from parent--child pairs.
\end{keyfinding}

\subsection{Formal Scaling Relations}
\label{sec:coverage:scaling}

Let $c$ denote coverage, $n(c)$ the module count, $m(c)$ the edge
count, and $\lambda_1(c)$ the spectral gap.  Under our null model:

\begin{align}
  n(c) &\approx N_0 \cdot \left(1 - e^{-\alpha c}\right)
  \label{eq:module_scaling} \\
  m(c) &\approx M_0 \cdot \left(1 - e^{-\beta c^2}\right)
  \label{eq:edge_scaling} \\
  \lambda_1(c) &\approx 1 - \sqrt{\frac{n(c)}{m(c)}} \cdot
  \lambda_1^{(\mathrm{reg})}
  \label{eq:gap_scaling}
\end{align}

where $N_0$ is the total number of possible modules, $M_0$ the total
number of possible edges, and $\alpha, \beta$ are detection sensitivity
parameters.  Equation~\eqref{eq:module_scaling} saturates as coverage
increases (all possible modules are eventually detected);
equation~\eqref{eq:edge_scaling} saturates faster due to the quadratic
dependence on coverage.

The implication is that \emph{any ecosystem feature that depends
monotonically on module count, edge count, or graph density will be
confounded by coverage}.  This includes most spectral features,
most graph statistics, and most straightforward comparisons.

\subsection{Proposed Deconfounding Strategies}
\label{sec:coverage:strategies}

We propose four strategies, ordered by increasing mathematical
sophistication.

\subsubsection{Strategy 1: Degree-Sequence Conditioning}

Compare the observed ecosystem to a \emph{null model} with the same
degree sequence.  Specifically, for each sample, generate $B$ random
graphs from the Chung--Lu model with the same expected degree sequence
as the observed ecosystem.  Compute the spectral features of each random
graph.  The \emph{deconfounded feature} is the z-score of the observed
feature relative to the null distribution:

\begin{equation}
  z(f) = \frac{f_{\mathrm{obs}} - \mathbb{E}[f_{\mathrm{null}}]}
              {\mathrm{Std}[f_{\mathrm{null}}]}
  \label{eq:zscore}
\end{equation}

This removes the effect of graph density (which is a function of
coverage) while preserving the effect of graph \emph{structure} (which
is a function of genomic organization).

\subsubsection{Strategy 2: Spectral Ratio Features}

Instead of using raw eigenvalues $\lambda_1, \lambda_2, \ldots$, use
\emph{ratios}: $\lambda_1 / \lambda_2$, $\lambda_2 / \lambda_3$, etc.
If coverage scales all eigenvalues by a common factor (which is
approximately true for dense graphs), ratios are invariant:

\begin{equation}
  \frac{\lambda_i(c)}{\lambda_j(c)} \approx
  \frac{\lambda_i}{\lambda_j} \quad \text{(coverage-independent)}
  \label{eq:ratio}
\end{equation}

This is the spectral analogue of using allele frequency \emph{ratios}
instead of raw counts in population genetics.

\subsubsection{Strategy 3: Rank-Based Comparison}

Convert all module features to ranks before computing ecosystem
similarity.  Ranks are invariant under any monotone transformation:
if coverage multiplies all features by a constant, their ranks are
unchanged.  Specifically, for each sample, rank the modules by each
feature dimension, then compare rank vectors between samples using
Spearman correlation.

\subsubsection{Strategy 4: Differential Subsampling}

For each pair of samples, progressively subsample reads (at 90\%, 80\%,
\ldots, 10\% of original coverage) and compute ecosystem similarity at
each level.  The \emph{slope} of the similarity-vs-coverage curve is
informative: related pairs should have a \emph{shallower slope}
(similarity decays slowly because the shared organizational structure
is robust) while unrelated pairs should have a \emph{steeper slope}
(similarity decays quickly because the background sharing is fragile).

\begin{equation}
  \Delta S = \frac{\partial S}{\partial c} \bigg|_{c = c_{\mathrm{obs}}}
  \approx \frac{S(0.9c) - S(0.5c)}{0.4c}
  \label{eq:slope}
\end{equation}

\begin{keyfinding}
\textbf{Differential subsampling converts a confound into a signal.}
Instead of trying to remove the coverage effect, we \emph{measure} the
coverage effect and use its magnitude as a feature.  Inherited structure
is robust to coverage reduction; coincidental similarity is not.
\end{keyfinding}


% ════════════════════════════════════════════════════════════════════════
% SECTION 6: CONNECTION TO CLASSICAL THEORY
% ════════════════════════════════════════════════════════════════════════
\section{Connection to Classical Theory}
\label{sec:classical}

\subsection{IBD as a Special Case of Ecosystem Correspondence}
\label{sec:classical:ibd}

Identity by descent (IBD) is the classical framework for genomic
relatedness~\cite{browning2012ibd}.  Two individuals share an IBD
segment if they inherited identical copies of a chromosomal region from
a common ancestor.  We show that IBD segments, viewed through the
ecosystem lens, correspond to subgraphs with \emph{perfect cooperative
stability}.

\begin{definition}
\textbf{IBD in the ecosystem framework.}  Let $S \subseteq V_A$ be the
set of modules in genome $A$ whose constituent $k$-mers fall within an
IBD segment shared with genome $B$, and let $S' \subseteq V_B$ be the
corresponding modules in $B$.  There exists a natural bijection
$\phi_S: S \to S'$ that maps each module to its IBD-identical
counterpart.  For this bijection:
\begin{enumerate}
  \item $\phi_S$ preserves all co-occurrence edges within the segment:
    $(m_i, m_j) \in E_A|_S \iff (\phi_S(m_i), \phi_S(m_j)) \in
    E_B|_{S'}$.
  \item $\phi_S$ preserves node features: $f_A(m_i) = f_B(\phi_S(m_i))$
    up to sequencing noise.
  \item The cooperative stability of $\phi_S$ restricted to $S$ is
    $S(\phi_S|_S) \approx 1$.
\end{enumerate}
\end{definition}

The coherence field $\Phi$ can be viewed as an \emph{approximate,
soft, global} version of this local bijection.  Over IBD segments,
$\Phi$ should approximate $\phi_S$.  Over non-IBD segments, $\Phi$ has
no template to follow and produces incoherent mappings.

\begin{pullquote}
IBD is a special case of ecosystem correspondence where the
correspondence is exact, local, and identity-based.  The ecosystem
framework generalizes IBD to \emph{approximate, global, structural}
correspondence.
\end{pullquote}

\subsection{Relationship to Linkage Disequilibrium}
\label{sec:classical:ld}

Linkage disequilibrium (LD) measures the non-random association of
alleles at different loci in a population.  Module co-occurrence in our
framework is analogous but distinct:

\begin{table}[H]
  \centering
  \tablestyle
  \caption{Comparison of linkage disequilibrium and module co-occurrence.}
  \label{tab:ld_comparison}
  \begin{tabular}{lll}
    \toprule
    \rowcolor{table-header}
    \textcolor{white}{\textbf{Property}} &
    \textcolor{white}{\textbf{LD}} &
    \textcolor{white}{\textbf{Module co-occurrence}} \\
    \midrule
    Level & Population (allele frequencies) & Individual (read-level) \\
    Reference & Allele at a locus & Module (cluster of $k$-mers) \\
    Scale & Chromosomal (cM) & Read-level (100--1000\,bp) \\
    Requires ref? & Yes (variant positions) & No (de novo modules) \\
    \bottomrule
  \end{tabular}
\end{table}

The ecosystem framework replaces population-level LD with
individual-level co-occurrence.  This has both advantages (no reference
genome, no population panel) and disadvantages (noisier, coverage-dependent).

\subsection{What the Ecosystem Approach Might Detect That IBD Cannot}
\label{sec:classical:beyond}

The strongest argument for the ecosystem approach is not that it will
beat IBD at detecting parent--child relationships (it probably cannot,
given the $100\times$ signal gap) but that it might detect phenomena
invisible to token-based IBD.

\begin{enumerate}
  \item \textbf{Ancestral organizational patterns.}  Organizational
    structure is more conserved than sequence identity.  Two individuals
    may share organizational patterns inherited from an ancestor many
    generations back, even though all specific IBD segments have been
    broken by recombination.

  \item \textbf{Population-level structural signatures.}  Different
    populations may have characteristic ecosystem structures (role
    distributions, spectral signatures) that persist across individuals.
    This could enable \emph{ancestry inference} from organizational
    structure.

  \item \textbf{Structural variants.}  Large structural variants
    (inversions, translocations, large deletions) disrupt $k$-mer
    identity across the breakpoint but may preserve the
    organizational structure within the rearranged segment.
    Liao et al.~\cite{liao2025sv} document extensive structural
    variation in human genomes; the ecosystem approach might be
    robust to these variants.

  \item \textbf{Shared regulatory grammar.}  Co-occurrence patterns
    among $k$-mers may reflect regulatory relationships (promoter
    motifs near coding motifs, repeat elements near gene boundaries).
    This ``regulatory grammar'' is not captured by haplotype identity
    but might be captured by ecosystem structure.

  \item \textbf{Metagenomics.}  In metagenomic samples (environmental
    DNA from microbial communities), there is no single reference genome.
    The ecosystem approach is inherently reference-free and could
    potentially be extended to compare community structure across
    samples.
\end{enumerate}


% ════════════════════════════════════════════════════════════════════════
% SECTION 7: EXPERIMENTAL STATUS
% ════════════════════════════════════════════════════════════════════════
\section{Experimental Status}
\label{sec:status}

\subsection{The CEPH 1463 Benchmark}
\label{sec:status:benchmark}

Our primary benchmark is the CEPH~1463 Platinum Pedigree, a
three-generation family sequenced by Illumina as part of the Platinum
Genomes project~\cite{ceph1463}.

\begin{pedigree}
\begin{verbatim}
                  GENERATION 1 (Grandparents)
    Pat.GF (NA12889)      Pat.GM (NA12890)
           \                 /
            \               /
    Mat.GF (NA12891)      Mat.GM (NA12892)
           \                 /
            \               /
                  GENERATION 2 (Parents)
            Father (NA12877) ---x--- Mother (NA12878)
                            |
                      GENERATION 3
                    Child (NA12882)
\end{verbatim}
\end{pedigree}

We use six samples (four grandparents plus father and mother) producing
$\binom{6}{2} = 15$ pairwise comparisons.  Of these:

\begin{itemize}
  \item \textbf{4 parent--child pairs}: Pat.GF$\to$Father,
    Pat.GM$\to$Father, Mat.GF$\to$Mother, Mat.GM$\to$Mother
  \item \textbf{1 spouse pair}: Father$\leftrightarrow$Mother
  \item \textbf{4 in-law pairs}: each grandparent with the opposite
    parent
  \item \textbf{6 cross-family pairs}: grandparents from different
    families
\end{itemize}

The benchmark is deliberately hard: all individuals are from the same
population (CEU/European), so the population sharing wall is
maximal.  Spouses provide a strong negative control: they share no
recent ancestry but were often sequenced with similar protocols and
depth.

\subsection{Current Leaderboard}
\label{sec:status:leaderboard}

Table~\ref{tab:leaderboard} presents the complete leaderboard of all
algorithm families tested on CEPH~1463.

\begin{table}[H]
  \centering
  \tablestyle
  \caption{Complete algorithm leaderboard on CEPH~1463.  Separation and
    gap as defined in Table~\ref{tab:token_methods}.  Status: PASS =
    positive gap (all parent--child above all non-parent--child);
    Fail = negative gap.}
  \label{tab:leaderboard}
  \begin{tabular}{llcccl}
    \toprule
    \rowcolor{table-header}
    \textcolor{white}{\textbf{Family}} &
    \textcolor{white}{\textbf{Best variant}} &
    \textcolor{white}{\textbf{Sep}} &
    \textcolor{white}{\textbf{Gap}} &
    \textcolor{white}{\textbf{Sp.}} &
    \textcolor{white}{\textbf{Status}} \\
    \midrule
    \multicolumn{6}{l}{\textit{Token-based methods}} \\
    Rare-run & P90 & +583\% & +300\% & low & \textbf{PASS} \\
    Multi-scale run & Weighted & +180\% & +50\% & low & \textbf{PASS} \\
    Run-length & Recombin.\ dist. & +12\% & $-$30\% & mid & Fail \\
    Distribution & Bray--Curtis & +3\% & $-$25\% & mid & Fail \\
    Information & NCD & +1.5\% & $-$40\% & mid & Fail \\
    Set overlap & Jaccard & $<$1\% & $-$50\% & high & Fail \\
    \midrule
    \multicolumn{6}{l}{\textit{Structural methods}} \\
    Module Persist. & v3 & +3.4\% & +0.13\% & low & \textbf{PASS} \\
    Coalition Trans. & v3 & +2.0\% & +0.06\% & low & \textbf{PASS} \\
    Module Stability & v4 & +4.9\% & $-$1.0\% & mid & Fail \\
    Spectral Ecol. & v7 & n/a & $-$0.6\% & mid & Fail \\
    Spectral Ecol. & v4 & +0.02\% & $-$5.0\% & \#8 & Fail \\
    Spectral Ecol. & v3 & +4.9\% & $-$54.7\% & \#1 & Fail \\
    \bottomrule
  \end{tabular}
\end{table}

\subsection{Infrastructure}
\label{sec:status:infrastructure}

The computational infrastructure comprises:

\begin{itemize}
  \item \textbf{12 TypeScript packages}: fastq, genome, kmer, alignment,
    relatedness, vcf, gpu, spirv, runtime, io, compression, math.
  \item \textbf{2 experiment packages}: symbiogenesis (module-level
    algorithms), ecology (spectral/graph algorithms).
  \item \textbf{18 ecology source files}: laplacian, heat-kernel,
    spectral-invariants, role-assignment, module-embedding,
    node-features, motif-graph, patch-tensor, coherence-field,
    multiscale-field, persistence-field, pair-transport, compare,
    ecosystem-signature, null-model, types, index, test suite.
  \item \textbf{4 symbiogenesis algorithms}: module-persistence,
    coalition-transfer, module-stability, plus barrel exports and
    runner.
  \item \textbf{LaTeX pipeline}: custom document class with wine palette,
    branded typography, build system targeting \texttt{tectonic}.
  \item \textbf{GPU pipeline} (stub): TypeScript DSL $\to$ SPIR-V IR
    $\to$ C codegen $\to$ GPU dispatch, with CPU fallbacks for all
    kernels.
  \item \textbf{${\sim}15{,}000$ lines} of TypeScript, all from scratch
    with no external bioinformatics libraries.
\end{itemize}

\subsection{Three Benchmark Datasets}
\label{sec:status:datasets}

\begin{table}[H]
  \centering
  \tablestyle
  \caption{Benchmark datasets used in our empirical program.}
  \label{tab:datasets}
  \begin{tabular}{llcl}
    \toprule
    \rowcolor{table-header}
    \textcolor{white}{\textbf{Dataset}} &
    \textcolor{white}{\textbf{Relationship}} &
    \textcolor{white}{\textbf{Samples}} &
    \textcolor{white}{\textbf{Difficulty}} \\
    \midrule
    Synthetic family & Parent-child trio & 3 & Easy \\
    GIAB Ashkenazi~\cite{zook2019giab} & Parent-child trio & 3 & Medium \\
    CEPH 1463~\cite{ceph1463} & 3-gen pedigree & 6 & Hard \\
    \bottomrule
  \end{tabular}
\end{table}

The synthetic family is useful for validating implementations but
trivial to separate.  The GIAB trio provides real sequencing noise.
CEPH~1463 is the definitive benchmark: same population, mixed
relationships, coverage variation, and the spouse confound.


% ════════════════════════════════════════════════════════════════════════
% SECTION 8: MATHEMATICAL DETAILS
% ════════════════════════════════════════════════════════════════════════
\section{Mathematical Details}
\label{sec:math}

This section provides the formal mathematical machinery underlying
our framework, intended for readers seeking rigorous foundations.

\subsection{The Population Sharing Probability}
\label{sec:math:sharing}

We derive the expected $k$-mer sharing between two individuals from the
same population.  Let $\pi$ denote the per-base SNP rate between two
random haplotypes.

For a $k$-mer to be \emph{identical} between two haplotypes, all $k$
positions must match.  Assuming SNPs are independently distributed
(a simplification that ignores LD):

\begin{equation}
  P(\text{$k$-mer match}) = (1 - \pi)^k
  \label{eq:kmer_match}
\end{equation}

For diploid organisms, each individual has two haplotypes.  A $k$-mer
is ``shared'' between individuals $A$ and $B$ if it appears on at
least one haplotype of each.  Let the four haplotype pairs be
$(h_A^1, h_B^1)$, $(h_A^1, h_B^2)$, $(h_A^2, h_B^1)$, $(h_A^2,
h_B^2)$.  By inclusion-exclusion:

\begin{align}
  P(\text{shared}) &= 1 - P(\text{not shared}) \nonumber \\
  &= 1 - P(\text{absent from $A$ or absent from $B$}) \nonumber \\
  &\approx 1 - (1 - (1-\pi)^k)^2 \cdot (1 - (1-\pi)^k)^2
  \label{eq:diploid_sharing}
\end{align}

At $\pi = 0.001$, $k = 21$:

\begin{align}
  (1 - \pi)^k &= 0.9792 \nonumber \\
  (1 - (1-\pi)^k)^2 &= (0.0208)^2 = 0.000433 \nonumber \\
  P(\text{shared}) &\approx 1 - 0.000433 \approx 0.9996
  \label{eq:sharing_numeric}
\end{align}

Virtually all 21-mers are shared.  The discriminative information
resides in the ${\sim}0.04\%$ of $k$-mers that are \emph{not}
universally shared.

\subsection{Expected Run Length Under IBD}
\label{sec:math:runlength}

Within an IBD segment, consecutive $k$-mers are shared unless a
sequencing error disrupts one.  The expected run length of shared
$k$-mers within an IBD segment of length $L$ bases is:

\begin{equation}
  \mathbb{E}[\text{run length} \mid \text{IBD}]
  \approx \frac{1}{\varepsilon_k} = \frac{1}{1 - (1-\varepsilon)^k}
  \label{eq:expected_run_ibd}
\end{equation}

where $\varepsilon$ is the per-base sequencing error rate.  At
$\varepsilon = 0.01$ and $k = 21$:

\begin{equation}
  \varepsilon_k = 1 - (1 - 0.01)^{21} = 1 - 0.810 = 0.190
  \label{eq:error_k}
\end{equation}

\begin{equation}
  \mathbb{E}[\text{run}] \approx \frac{1}{0.190} \approx 5.3
  \label{eq:expected_run_numeric}
\end{equation}

Outside IBD, the expected run length is much shorter, determined by
the chance overlap probability at the rare-$k$-mer level.  The run-length
distribution is the key discriminative feature exploited by the
rare-run P90 method.

\subsection{Module Count as a Function of Coverage}
\label{sec:math:modules}

Let $M$ be the total number of distinct modules in the genome (a
property of the genome, not the sequencing).  At coverage $c$, a module
with population frequency $p$ and $s$ constituent $k$-mers is
\emph{detected} if at least $\tau$ of its $k$-mers are observed.  The
number of observed $k$-mers follows $\mathrm{Binomial}(s, 1 -
(1 - p)^c) \approx \mathrm{Binomial}(s, 1 - e^{-pc})$ for large $c$.

The expected number of detected modules is:

\begin{equation}
  n(c) = \sum_{m=1}^{M} P\!\left(
    \mathrm{Bin}(s_m, 1 - e^{-p_m c}) \geq \tau
  \right)
  \label{eq:module_count}
\end{equation}

For typical parameters ($p_m \approx 0.5$, $s_m \approx 10$,
$\tau = 3$), this saturates at $c \approx 10\times$, well within
standard whole-genome sequencing depth.

\subsection{Gromov--Wasserstein: Entropic Regularization}
\label{sec:math:gw_entropic}

The exact GW problem (Equation~\ref{eq:gw}) is computationally
intractable for large $n$.  The entropic-regularized version
adds a penalty $\varepsilon \, H(\pi)$ where $H(\pi) = -\sum_{ij}
\pi_{ij} \log \pi_{ij}$ is the entropy of the transport plan:

\begin{equation}
  \mathrm{GW}_\varepsilon(X, Y) = \inf_{\pi \in \Pi(\mu_X, \mu_Y)}
  \left[
    \sum_{i,i'} \sum_{j,j'}
    |d_X(x_i, x_{i'}) - d_Y(y_j, y_{j'})|^2 \,
    \pi_{ij} \pi_{i'j'}
    - \varepsilon \, H(\pi)
  \right]
  \label{eq:gw_entropic}
\end{equation}

This can be solved by alternating between:
\begin{enumerate}
  \item Fixing $\pi$ and solving the linear optimal transport problem
    (Sinkhorn iterations~\cite{cuturi2013sinkhorn}).
  \item Fixing the cost matrix and updating $\pi$ via a conditional
    gradient step.
\end{enumerate}

Each Sinkhorn step is $O(n^2)$; typically 10--50 outer iterations
suffice for convergence.  Total cost: $O(n^2 \cdot K_{\mathrm{Sinkhorn}}
\cdot K_{\mathrm{outer}})$ where $K_{\mathrm{Sinkhorn}} \approx 50$ and
$K_{\mathrm{outer}} \approx 20$.  For $n = 700$: approximately
$700^2 \times 50 \times 20 = 4.9 \times 10^8$ operations, feasible in
seconds on a modern CPU.

\subsection{Persistence Diagram Stability}
\label{sec:math:stability}

The Bottleneck Stability Theorem guarantees that small perturbations of
the input produce small changes in the persistence diagram.

\begin{definition}
\textbf{Bottleneck Stability Theorem}
(Edelsbrunner et al.~\cite{edelsbrunner2010computational}).  Let $f$
and $g$ be two tame functions on a topological space $X$, and let
$\mathrm{Dgm}(f)$ and $\mathrm{Dgm}(g)$ denote their persistence
diagrams.  Then:
\begin{equation}
  d_B(\mathrm{Dgm}(f), \mathrm{Dgm}(g)) \leq \|f - g\|_\infty
  \label{eq:bottleneck_stability}
\end{equation}
where $d_B$ is the bottleneck distance between diagrams.
\end{definition}

For genomic ecosystems, this means: if two samples have similar
co-occurrence weight functions (because they share inherited
organizational structure), their persistence diagrams will be close in
bottleneck distance.  Sequencing noise, which adds small perturbations
to edge weights, produces correspondingly small changes in the diagram.

\subsection{Connection Curvature on Graphs}
\label{sec:math:connection}

We formalize the gauge-theoretic direction.  Let $\mathcal{G} = (V, E)$
be a graph with a vector field $\Phi: V \to \mathbb{R}^d$ (the
coherence field).  Define the \emph{discrete connection} $\omega$ along
an edge $(i, j) \in E$ as:

\begin{equation}
  \omega_{ij} = \Phi(j) - \Phi(i)
  \label{eq:discrete_connection}
\end{equation}

The \emph{discrete curvature} around a triangle $(i, j, k)$ is the
holonomy:

\begin{equation}
  \Omega_{ijk} = \omega_{ij} + \omega_{jk} + \omega_{ki}
  = [\Phi(j) - \Phi(i)] + [\Phi(k) - \Phi(j)] + [\Phi(i) - \Phi(k)]
  = 0
  \label{eq:zero_curvature}
\end{equation}

This is identically zero!  The trivial connection has zero curvature.
To obtain a non-trivial curvature, we must introduce a non-trivial
parallel transport.  The natural choice is:

\begin{equation}
  \omega_{ij}^{\mathrm{OT}} = \Phi(j) - T_{ij}(\Phi(i))
  \label{eq:ot_connection}
\end{equation}

where $T_{ij}$ is the \emph{optimal transport map} from the local
neighbourhood of $i$ (in $B$-space) to the local neighbourhood of $j$
(in $B$-space).  The curvature then measures whether going around
a loop in $A$'s graph and applying the local transport maps at each
step returns to the starting point in $B$'s space.

\begin{definition}
\textbf{Holonomy curvature.}  For a triangle $(i, j, k) \in
T(\mathcal{G}_A)$:
\begin{equation}
  \Omega_{ijk}^{\mathrm{OT}} = \omega_{ij}^{\mathrm{OT}}
  + \omega_{jk}^{\mathrm{OT}} + \omega_{ki}^{\mathrm{OT}}
  \label{eq:holonomy}
\end{equation}
A ``flat'' field ($\Omega \approx 0$) indicates that the local transport
maps are globally consistent---the structural correspondence between $A$
and $B$ is coherent.  A ``curved'' field ($\|\Omega\|$ large) indicates
global inconsistency.
\end{definition}

This is directly analogous to the curvature of a connection on a
principal bundle in differential geometry~\cite{nakahara2003geometry},
discretized to the graph setting.


% ════════════════════════════════════════════════════════════════════════
% SECTION 9: OPEN QUESTIONS
% ════════════════════════════════════════════════════════════════════════
\section{Open Questions}
\label{sec:openquestions}

We close by identifying the most important open questions, both
theoretical and empirical, that this research program must address.

\subsection{Is the Structural Signal Strong Enough?}
\label{sec:open:strength}

The central empirical question: is the organizational structure of a
genome sufficiently different between related and unrelated individuals
to be practically useful?  Our current results suggest the signal
exists ($+3.36\%$ separation for Module Persistence v3, positive
weakest gap) but is $100\times$ weaker than the token-level signal
($+583\%$ for rare-run P90).

There are two possible interpretations:

\begin{enumerate}
  \item \textbf{The signal is weak because our methods are immature.}
    Better mathematical tools (GW distance, persistence, gauge
    curvature) might extract orders of magnitude more structural
    signal.
  \item \textbf{The signal is inherently weak because organizational
    structure varies less than sequence identity.}  Related and
    unrelated individuals may simply have more similar organizational
    structures than they have similar sequence content.
\end{enumerate}

Distinguishing these hypotheses is critical.  If (2) is correct, the
ecosystem approach has fundamental limits for close-relationship
detection.  It might still be valuable for distant-relationship
detection, ancestry inference, or metagenomics, where token-level
signal is insufficient.

\subsection{Can Spectral Methods Match Rare-Run Performance?}
\label{sec:open:match}

Specifically for parent--child detection on CEPH~1463: can any spectral
or structural method achieve $+50\%$ separation (currently the
second-best token result)?  Or is the spectral approach better suited
for different tasks?

We conjecture that spectral methods will never match rare-run P90 at
parent--child detection, because rare-run is essentially doing IBD
detection and IBD is the \emph{strongest possible signal} for close
relationships.  The spectral approach should instead target:

\begin{itemize}
  \item Distant relationships (beyond the IBD detection horizon)
  \item Population structure (where organizational patterns may be
    shared across many individuals)
  \item Metagenomic community structure (where no reference exists)
\end{itemize}

\subsection{Does the Ecosystem Representation Contain Novel Information?}
\label{sec:open:novel}

A critical theoretical question: does the ecosystem graph $\mathcal{G}$
encode information that is \emph{not} present in the $k$-mer frequency
spectrum $\{(k_i, c_i)\}$ (the multiset of $k$-mers and their counts)?

We argue yes: the ecosystem graph encodes \emph{co-occurrence structure}
(which $k$-mers appear together in reads) which is not captured by the
marginal frequency spectrum.  Two genomes can have identical $k$-mer
frequency spectra but different co-occurrence graphs if the $k$-mers
are arranged differently along the chromosomes.

However, this argument depends on read length.  With infinitely long
reads, the co-occurrence graph would capture the entire chromosomal
arrangement and contain maximal information.  With short reads
(${\sim}150$\,bp for Illumina), the co-occurrence graph captures only
\emph{local} arrangement within a $150$\,bp window.  The
information content is therefore bounded by the ratio of read length to
genome size.

\subsection{Information-Theoretic Limits}
\label{sec:open:info}

What is the information-theoretic limit of reference-free relatedness
detection from raw reads?  Formally: given $N$ reads of length $L$ from
a genome of size $G$ at coverage $c = NL/G$, what is the mutual
information between the read set and the relatedness status
(related/unrelated)?

This is an open theoretical question.  A lower bound is provided by
the token-level information (the rare-run P90 signal).  An upper bound
is provided by the total information in the read set about the donor
genome.  The gap between these bounds determines how much room exists
for structural methods to improve upon token methods.

\subsection{Extension to Metagenomics}
\label{sec:open:meta}

The ecosystem framework is inherently reference-free and operates on
communities of interacting motifs.  This makes it conceptually
applicable to metagenomics, where the ``genome'' is a community of
organisms and the ``modules'' are shared gene families or metabolic
pathways.

The open question is whether the spectral-geometric tools developed
for single-organism relatedness detection transfer to community-level
comparison.  The scale is different (millions of species vs.\ hundreds
of modules), but the mathematical framework---interaction graphs,
Laplacian spectra, coherence fields---is the same.


% ════════════════════════════════════════════════════════════════════════
% SECTION 10: CONCLUSION
% ════════════════════════════════════════════════════════════════════════
\section{Conclusion}
\label{sec:conclusion}

\subsection{An Honest Assessment}

The token paradigm works.  On the hardest benchmark we have (CEPH~1463,
same-population, three-generation pedigree), the rare-run P90 achieves
$+583\%$ separation with every parent--child pair scoring decisively
above every non-parent--child pair.  This is a practical, deployable
result that validates reference-free relatedness detection as an
engineering possibility.

The ecosystem paradigm is promising but unproven.  Module Persistence v3
and Coalition Transfer v3 achieve positive weakest gap on CEPH~1463,
demonstrating that organizational structure carries kinship information.
But the signal is $100\times$ weaker than the token signal, and the
spectral ecology program---despite seven iterations of increasing
mathematical sophistication---has not yet achieved positive weakest gap.

\subsection{What the Trajectory Demonstrates}

The spectral ecology arc (v1 through v7) demonstrates several important
facts:

\begin{enumerate}
  \item \textbf{Structural signal exists.}  The improvements from
    $-26.87\%$ to $-0.61\%$ are not random; each version targets and
    addresses a specific failure mode.

  \item \textbf{The spouse confound is breakable.}  v4's breakthrough
    (spouses from rank~1 to rank~8) proves that structural methods can
    distinguish ``similar because co-sequenced'' from ``similar because
    related''---something naive token methods cannot do.

  \item \textbf{Diminishing returns within the kernel-regression
    paradigm.}  The $44\times$ improvement from v1 to v7 is large,
    but the rate of improvement is decelerating.  Qualitative
    mathematical leaps (GW, persistence, gauge theory) are needed.

  \item \textbf{The computational objects are converging.}  The
    progression tokens $\to$ summaries $\to$ matchings $\to$ fields
    $\to$ multi-scale fields $\to$ persistence fields suggests a
    natural hierarchy.  Each level adds more structural information
    to the comparison.
\end{enumerate}

\subsection{The Path Forward}

We believe the strongest case for the ecosystem approach will come
not from beating rare-run P90 at parent--child detection, but from
one of three scenarios:

\begin{enumerate}
  \item \textbf{Detecting distant relationships} beyond the IBD
    detection horizon, where organizational patterns persist after
    specific sequence identity has been lost to recombination.

  \item \textbf{Detecting population structure} from organizational
    signatures, enabling reference-free ancestry inference without
    variant calling.

  \item \textbf{Detecting structural relationships in metagenomics},
    where no reference genome exists and the ecosystem framework
    is the only available paradigm.
\end{enumerate}

The three proposed theoretical directions---Gromov--Wasserstein distance,
persistent homology, and gauge-theoretic field comparison---represent
our best mathematical candidates for qualitative advances.  Each
addresses a specific limitation of the current coherence field framework
(size dependence, spectral blindness, and outlier sensitivity,
respectively), and each draws on mature mathematical theory with strong
stability guarantees.

\begin{pullquote}
The question is not whether genomes can be compared structurally---they
can.  The question is whether structural comparison reveals anything
that sequence comparison does not.
\end{pullquote}

This paper is a request for feedback.  The theoretical framework, the
empirical trajectory, and the proposed directions are presented in full
so that computational biologists, applied mathematicians, and genomicists
can evaluate the approach, identify flaws, and suggest alternatives.
We welcome criticism.  The strongest possible outcome for this program
is a rigorous answer---whether positive or negative---to the question of
whether inherited ecosystem structure is a practically useful signal
for genomic relatedness detection.


% ════════════════════════════════════════════════════════════════════════
% APPENDIX: ALGORITHMIC CATALOGUE
% ════════════════════════════════════════════════════════════════════════
\appendix
\section{Algorithmic Catalogue}
\label{sec:catalogue}

This appendix provides brief descriptions of all major algorithm
variants tested during the empirical program, organized by family.

\subsection{Set Overlap Family}

\begin{table}[H]
  \centering
  \tablestyle
  \caption{Set overlap algorithms.}
  \label{tab:cat_set}
  \begin{tabular}{llc}
    \toprule
    \rowcolor{table-header}
    \textcolor{white}{\textbf{Algorithm}} &
    \textcolor{white}{\textbf{Description}} &
    \textcolor{white}{\textbf{CEPH Sep}} \\
    \midrule
    $k$-mer Jaccard & $|A \cap B| / |A \cup B|$ & $<1\%$ \\
    Containment & $|A \cap B| / \min(|A|, |B|)$ & $<1\%$ \\
    MinHash & Approximate Jaccard via sketching & $<1\%$ \\
    Weighted Jaccard & IDF-weighted set intersection & $\sim$2\% \\
    \bottomrule
  \end{tabular}
\end{table}

All set overlap methods fail because the population sharing wall
($\sim$98\% baseline) dominates the IBD signal ($\sim$1\% excess).

\subsection{Run-Length Family}

\begin{table}[H]
  \centering
  \tablestyle
  \caption{Run-length algorithms.}
  \label{tab:cat_run}
  \begin{tabular}{llc}
    \toprule
    \rowcolor{table-header}
    \textcolor{white}{\textbf{Algorithm}} &
    \textcolor{white}{\textbf{Description}} &
    \textcolor{white}{\textbf{CEPH Sep}} \\
    \midrule
    Recombination distance & Mean run length of shared $k$-mers & +12\% \\
    P90 run length & 90th percentile of run lengths & +25\% \\
    Weighted runs & Length-weighted run score & +30\% \\
    Multi-scale runs & Runs at multiple $k$ values & +180\% \\
    Rare P90 & P90 on rare ($<$50\% freq) $k$-mers & \textbf{+583\%} \\
    \bottomrule
  \end{tabular}
\end{table}

The rare-run P90 breakthrough demonstrates that rarity filtering
removes the population baseline, leaving only the IBD-informative
$k$-mers.

\subsection{Distribution Family}

\begin{table}[H]
  \centering
  \tablestyle
  \caption{Distribution comparison algorithms.}
  \label{tab:cat_dist}
  \begin{tabular}{llc}
    \toprule
    \rowcolor{table-header}
    \textcolor{white}{\textbf{Algorithm}} &
    \textcolor{white}{\textbf{Description}} &
    \textcolor{white}{\textbf{CEPH Sep}} \\
    \midrule
    Cosine similarity & Frequency vector cosine & +2\% \\
    Jensen--Shannon & $\sqrt{\mathrm{JSD}(p_A, p_B)}$ & +1.5\% \\
    Bray--Curtis~\cite{bray1957} & $1 - \sum|a_i - b_i| / \sum(a_i + b_i)$ & +3\% \\
    Morisita--Horn~\cite{morisita1959} & Overlap index for count data & +2.5\% \\
    \bottomrule
  \end{tabular}
\end{table}

Distribution methods provide moderate signal but cannot distinguish
local sharing (IBD) from global sharing (population baseline).

\subsection{Information-Theoretic Family}

\begin{table}[H]
  \centering
  \tablestyle
  \caption{Information-theoretic algorithms.}
  \label{tab:cat_info}
  \begin{tabular}{llc}
    \toprule
    \rowcolor{table-header}
    \textcolor{white}{\textbf{Algorithm}} &
    \textcolor{white}{\textbf{Description}} &
    \textcolor{white}{\textbf{CEPH Sep}} \\
    \midrule
    Entropy rate & Conditional entropy of shared $k$-mers & +1\% \\
    NCD~\cite{li2004ncd} & Normalized compression distance & +1.5\% \\
    Mutual information & MI between $k$-mer count vectors & +0.5\% \\
    \bottomrule
  \end{tabular}
\end{table}

Information-theoretic methods perform poorly, likely because the
compressibility of genomic sequences is dominated by repeat content
rather than by kinship-related sharing patterns.

\subsection{Structural: Module Persistence Family}

Module Persistence constructs a shared module vocabulary across all
samples and compares samples by which modules they contain.

\begin{itemize}
  \item \textbf{v1}: Per-sample modules with Jaccard comparison.
    Fails because module IDs are sample-specific.
  \item \textbf{v2}: Shared vocabulary with edge Jaccard.  Edge Jaccard
    has range 86--89\% for all pairs (not discriminative).
  \item \textbf{v3}: Rarity-weighted informative module overlap plus
    support cosine.  \textbf{$+3.36\%$ separation, $+0.13\%$ gap.  PASSES.}
\end{itemize}

\subsection{Structural: Coalition Transfer Family}

Coalition Transfer measures sharing of \emph{module coalitions}:
specific combinations of modules that co-occur in individual reads.

\begin{itemize}
  \item \textbf{v1}: Raw coalition Jaccard.  Moderate signal but
    confounded by coverage.
  \item \textbf{v2}: IDF-weighted coalition cosine.  Boosted in-laws
    unexpectedly.
  \item \textbf{v3}: Three-way blend (rarity + exclusive + unweighted
    cosine).  \textbf{$+1.97\%$ separation, $+0.06\%$ gap.  PASSES.}
\end{itemize}

\subsection{Structural: Module Stability Family}

Module Stability uses bootstrap resampling to identify which motifs
are \emph{stably organized} into modules across subsamples.

\begin{itemize}
  \item \textbf{v1--v3}: Binary stability filtering with various
    thresholds.  Moderate signal.
  \item \textbf{v4}: Continuous stability-weighted rarity Jaccard.
    \textbf{$+4.90\%$ average separation} but \textbf{$-1.04\%$ weakest
    gap} (fails weakest parent--child pair).
\end{itemize}

\subsection{Structural: Spectral Ecology Family}

The Spectral Ecology family is detailed in
Section~\ref{sec:trajectory:arc} and represents our deepest
investigation of spectral-geometric methods.  Seven versions span the
progression from disconnected sequential-adjacency graphs (v1) through
Sinkhorn-regularized optimal transport with cooperative stability
weighting (v7).


% ════════════════════════════════════════════════════════════════════════
% APPENDIX: CEPH 1463 PAIRWISE ANALYSIS
% ════════════════════════════════════════════════════════════════════════
\section{CEPH 1463 Detailed Pairwise Analysis}
\label{sec:pairwise}

For completeness, we present the complete pairwise score tables for
the three best-performing algorithms on CEPH~1463.

\subsection{Module Persistence v3}

\begin{table}[H]
  \centering
  \tablestyle
  \caption{Module Persistence v3: complete pairwise scores on CEPH~1463.
    $\blacktriangleright$ = parent--child.}
  \label{tab:mp_scores}
  \begin{tabular}{clccc}
    \toprule
    \rowcolor{table-header}
    \textcolor{white}{\textbf{Rk}} &
    \textcolor{white}{\textbf{Pair}} &
    \textcolor{white}{\textbf{Score}} &
    \textcolor{white}{\textbf{Rarity}} &
    \textcolor{white}{\textbf{Cosine}} \\
    \midrule
    1 & $\blacktriangleright$ Pat.GM $\to$ Father & 0.3248 & 34.1\% & 0.320 \\
    2 & $\blacktriangleright$ Mat.GF $\to$ Mother & 0.3221 & 33.8\% & 0.318 \\
    3 & $\blacktriangleright$ Pat.GF $\to$ Father & 0.3195 & 33.4\% & 0.316 \\
    4 & $\blacktriangleright$ Mat.GM $\to$ Mother & 0.3148 & 32.9\% & 0.312 \\
    \midrule
    5 & Father $\leftrightarrow$ Mother & 0.3144 & 32.8\% & 0.311 \\
    6 & Mat.GF $\leftrightarrow$ Father & 0.3132 & 32.7\% & 0.310 \\
    7 & Pat.GF $\leftrightarrow$ Mother & 0.3118 & 32.5\% & 0.309 \\
    8 & Pat.GM $\leftrightarrow$ Mat.GM & 0.3105 & 32.3\% & 0.308 \\
    9 & Pat.GF $\leftrightarrow$ Mat.GF & 0.3098 & 32.2\% & 0.307 \\
    10 & Mat.GM $\leftrightarrow$ Father & 0.3091 & 32.1\% & 0.307 \\
    \bottomrule
  \end{tabular}
\end{table}

\ymarginnote{All four parent--child pairs rank in the top four.  Gap:
$+0.13\%$.  Razor-thin but correct.}

The margins are razor-thin: the gap between rank~4 (weakest
parent--child: 0.3148) and rank~5 (strongest non-parent--child: 0.3144)
is only $+0.13\%$.  This demonstrates that the structural signal is
real but fragile.

\subsection{Coalition Transfer v3}

\begin{table}[H]
  \centering
  \tablestyle
  \caption{Coalition Transfer v3: complete pairwise scores on
    CEPH~1463.}
  \label{tab:ct_scores}
  \begin{tabular}{clccc}
    \toprule
    \rowcolor{table-header}
    \textcolor{white}{\textbf{Rk}} &
    \textcolor{white}{\textbf{Pair}} &
    \textcolor{white}{\textbf{Score}} &
    \textcolor{white}{\textbf{Excl.}} &
    \textcolor{white}{\textbf{FreqCos}} \\
    \midrule
    1 & $\blacktriangleright$ Pat.GF $\to$ Father & 0.1524 & 1.82\% & 0.394 \\
    2 & $\blacktriangleright$ Pat.GM $\to$ Father & 0.1498 & 1.75\% & 0.391 \\
    3 & $\blacktriangleright$ Mat.GF $\to$ Mother & 0.1472 & 1.68\% & 0.387 \\
    4 & $\blacktriangleright$ Mat.GM $\to$ Mother & 0.1463 & 1.65\% & 0.385 \\
    \midrule
    5 & Father $\leftrightarrow$ Mother & 0.1462 & 1.64\% & 0.384 \\
    6 & Mat.GF $\leftrightarrow$ Father & 0.1455 & 1.62\% & 0.383 \\
    7 & Pat.GF $\leftrightarrow$ Mother & 0.1448 & 1.59\% & 0.381 \\
    8 & Pat.GM $\leftrightarrow$ Mat.GM & 0.1441 & 1.57\% & 0.380 \\
    9 & Pat.GF $\leftrightarrow$ Mat.GF & 0.1438 & 1.56\% & 0.379 \\
    10 & Mat.GM $\leftrightarrow$ Father & 0.1432 & 1.54\% & 0.378 \\
    \bottomrule
  \end{tabular}
\end{table}

Coalition Transfer achieves the same perfect ordinal ranking as Module
Persistence but with even thinner margins ($+0.06\%$ gap).


% ════════════════════════════════════════════════════════════════════════
% APPENDIX: NOTATION AND SYMBOLS
% ════════════════════════════════════════════════════════════════════════
\section{Notation and Symbols}
\label{sec:notation}

\begin{table}[H]
  \centering
  \tablestyle
  \caption{Notation used throughout this paper.}
  \label{tab:notation}
  \begin{tabular}{ll}
    \toprule
    \rowcolor{table-header}
    \textcolor{white}{\textbf{Symbol}} &
    \textcolor{white}{\textbf{Meaning}} \\
    \midrule
    $\mathcal{G} = (V, E, w, f)$ & Genomic ecosystem graph \\
    $L_{\mathrm{sym}}$ & Symmetric normalized Laplacian \\
    $\lambda_k$ & $k$-th eigenvalue of $L_{\mathrm{sym}}$ \\
    $\phi_k$ & $k$-th eigenvector of $L_{\mathrm{sym}}$ \\
    $K_t = \exp(-tL)$ & Heat kernel at timescale $t$ \\
    $\mathrm{HKS}(v, t)$ & Heat kernel signature of node $v$ \\
    $\Phi: \mathcal{G}_A \to \mathbb{R}^d$ & Coherence field \\
    $S(\Phi)$ & Cooperative stability of field $\Phi$ \\
    $\mathrm{GW}(X, Y)$ & Gromov--Wasserstein distance \\
    $\pi \in \Pi(\mu, \nu)$ & Optimal transport plan \\
    $H(\pi)$ & Entropy of transport plan \\
    $D_A, D_B$ & Persistence diagrams \\
    $W_p(D_A, D_B)$ & $p$-Wasserstein distance between diagrams \\
    $\omega_{ij}$ & Discrete connection along edge $(i,j)$ \\
    $\Omega_{ijk}$ & Holonomy curvature around triangle \\
    $\mathcal{R}(\Phi)$ & Aggregate curvature of field \\
    $c$ & Sequencing coverage \\
    $k$ & $k$-mer length \\
    $\pi$ & Per-base SNP rate \\
    $n$ & Number of modules \\
    $m$ & Number of edges \\
    $d$ & Patch tensor dimension (13) \\
    $\sigma$ & Kernel bandwidth \\
    \bottomrule
  \end{tabular}
\end{table}


% ════════════════════════════════════════════════════════════════════════
% BIBLIOGRAPHY
% ════════════════════════════════════════════════════════════════════════
\bibliographystyle{plain}
\bibliography{../references}

\end{document}
